% v005 Related Work, assembled 2026-11-17. Kept verbatim from v003: "Evolutionary NAS with surrogates",
% the pyramid figure and "Joint NAS+HPO benchmarks". New: positioning.tex (replaces v003 "Linkage learning
% and P3", "Surrogate-assisted linkage learning", "Scale dependence", "Closest prior work"),
% compared_methods.tex, llm.tex (replaces the v003 LLM paragraph).

\section{Related Work}
\label{sec:related-work}

\paragraph{Evolutionary NAS with surrogates.} NSGA-Net searches architectures with NSGA-II, training every
candidate fully~\citep{lu2019nsganet}. NSGANetV2 adds an accuracy predictor refined online and inherits weights
from a supernet, cutting the number of fully evaluated architectures to 350~\citep{lu2020nsganetv2}, at the
documented risk of rank disagreement between weight-shared and stand-alone accuracy~\citep{yu2020evaluatingnas}.
White et al.\ identify multi-objective evolutionary search as the most common approach to multi-objective NAS,
with NSGANetV2 as an example~\citep{white2023insights}; Liu et al.\ list learned surrogates as one of several
cost-reduction techniques alongside weight sharing and zero-cost proxies~\citep{liu2023survey}. Later work in
this lineage moves toward ranking-based predictors~\citep{lu2022mosegnas}. None of three recent surveys lists
linkage learning as a distinct NAS method category~\citep{liu2023survey,white2023insights,ozcelik2026survey}.

% v005 stage C draft (2026-11-17): positioning of P3Net in the P3/GOMEA line. Replaces the v003 paragraphs
% "Linkage learning and P3", "Surrogate-assisted linkage learning" and "Closest prior work"
% (chapters/archive/v003/related_work/main.tex); the v003 pyramid figure (fig:p3-pyramid) can be reused unchanged.
% Sources: notes/ga_landscape_2026.md §16-§16d (literature search 2026-11-16/17). Dahi et al. 2026 is described
% from its supplementary material only (main text not accessible). The sentence on the surrogate result follows the
% current v005 + E1 results and is to be revised if the surrogate study continues.

\paragraph{Linkage learning and P3.} GOMEA and P3 build a linkage tree from statistical dependencies in the
population and use it to guide crossover, so that groups of dependent variables are not
disrupted~\citep{thierens2011optimal,goldman2014parameterless}. P3 replaces a single fixed-size population with a
pyramid of populations (Figure~\ref{fig:p3-pyramid}), adding a level only once existing levels stop improving, so that no population size is
chosen in advance; its evaluation complexity grows more slowly with problem size than that of five
competitors~\citep{goldman2015fast}. Multi-objective variants exist for both: MO-GOMEA~\citep{luong2014mogomea}
and MO-P3, applied to industrial process planning~\citep{przewozniczek2021mop3} and later extended with a Pareto-front
improvement phase~\citep{niemczyk2025pareto}. Unlike the gray-box setting that this line often
exploits~\citep{whitley2016graybox,bouter2023library}, NAS and HPO objectives are not decomposable into cheaply
recomputable subfunctions, so every variation step must be verified by a full evaluation or screened by a surrogate.

\paragraph{Linkage learning in NAS.} Members of this family have been applied to NAS before. den Ottelander et al.\
find MO-GOMEA, NSGA-II, local search and random search indistinguishable on NAS-Bench-101, with MO-GOMEA ahead of
random search on their larger MacroNAS spaces~\citep{denottelander2021localsearch}. ENCAS uses MO-GOMEA to search
cascades of networks drawn from supernetworks~\citep{chebykin2022encas}; Tran et al.\ apply GOMEA to NAS with a
training-free metric in place of training~\citep{tran2023gomeanas}; and Dahi et al., after applying linkage-learning
EAs to architecture search for quantum-circuit transpilation~\citep{dahi2024evolutionary}, use P3 to choose the depth
and layer widths of a small multilayer perceptron within a quantum-optimisation pipeline, a single-objective
architecture search~\citep{dahi2026tightening}. Bartnik adapts the surrogate-assisted P3-GOMEA
of~\citet{dushatskiy2021novelsurrogateassistedevolutionaryalgorithm} to multi-objective NAS on NAS-Bench-201
(accuracy against GPU energy), with an absolute regressor over the genotype as surrogate, an architecture-only
genotype, and MO-LS and surrogate-free MO-P3-GOMEA as baselines~\citep{bartnik2026evolutionary}; we include its two
variants, MO-DSA and MO-PSA, among our comparison arms.

\paragraph{Surrogates and linkage.} CS-GOMEA pairs GOMEA with a convolutional surrogate and names as open the
surrogate's failure to use linkage information in its own construction~\citep{dushatskiy2019csgomea}. LyMPuS builds
the surrogate from the discovered linkage itself and validates it with P3 on NK landscapes, Ising spin glasses and
Max3Sat; its empirical variant, eLyMPuS, discovers the dependency structure incrementally rather than assuming it
known~\citep{przewozniczek2026lympus}, and related work constructs monotonic surrogates from recursive linkage
discovery~\citep{przewozniczek2026monotonical}. P3Net's surrogate follows eLyMPuS in being relative and
linkage-aware over an incompletely discovered structure, but regresses a continuous fitness difference between
two solutions instead of making a discrete better/worse/ambiguous comparison, which places it closer to CS-GOMEA's
learned model and gives up eLyMPuS's guarantee of recovering the comparison under monotonicity.

\paragraph{Position of this study.} Against this background, P3Net is not the first use of linkage learning, or of
P3, in architecture search. What we have not found is P3 applied to the \emph{joint} multi-objective search over
architecture and training hyperparameters, nor a comparison of a linkage-learning optimiser with the main
Bayesian-optimisation, multi-fidelity and evolutionary methods of multi-objective NAS and HPO under one
pre-registered protocol that also reports the optimiser's own cost. This study supplies that comparison, with P3
in three forms (canonical, with a pairwise surrogate, and with eLyMPuS), on four primary spaces and five further
ones, and it reports where the surrogate does and does not change the outcome. All benchmarks are tabular or
surrogate-based (Section~\ref{sec:setup}).


\begin{figure}[t]
    \centering
    \begin{tikzpicture}[
        lvl/.style={draw, rounded corners, minimum height=6mm, font=\scriptsize, align=center},
    ]
        \node[lvl, minimum width=14mm] (l1) at (0,0)     {level 1};
        \node[lvl, minimum width=22mm] (l2) at (0,-9mm)  {level 2};
        \node[lvl, minimum width=30mm] (l3) at (0,-18mm) {level 3};
        \node[lvl, minimum width=38mm] (l4) at (0,-27mm) {level 4};
        \draw[{Stealth[length=2mm]}-, dashed] (l4.south) -- ++(0,-6mm)
            node[below, font=\tiny, align=center, text width=32mm]
            {new level added only once existing levels stop yielding improved solutions};
    \end{tikzpicture}
    \caption{P3's population pyramid: each level holds a population of growing size; a new level is added only
    when the levels below it stop producing improved solutions.}
    \Description{A vertical stack of four boxes labelled level 1 through level 4, growing wider from top to
    bottom, with an arrow noting that a new level is added only once existing levels stop yielding improved
    solutions.}
    \label{fig:p3-pyramid}
\end{figure}

% v005 stage C draft (2026-11-17): paragraph on the compared methods, to join the v003 Related Work section.
% Citations only for methods actually run in v005 (configs/pipeline/v005.yaml, v005-vizier.yaml). The P3/GOMEA
% positioning paragraphs wait for the author's literature search (notes/archive/v005/plans/v005-todo.md, stage C).

\paragraph{Compared methods.} The comparison covers the main families of multi-objective NAS and HPO optimisers,
each in its published form or the closest available implementation. \emph{NAS-specific search:} NSGA-Net and
NSGANetV2~\citep{lu2019nsganet,lu2020nsganetv2}, REINFORCE~\citep{williams1992reinforce,zoph2017nas} in the
NAS-Bench-201 reference implementation~\citep{dong2020nasbench201} with the MONAS reward~\citep{hsu2018monas} ($\alpha = 0.75$), and regularised evolution~\citep{real2019regularized}.
\emph{Linkage-learning and local search:} canonical P3~\citep{goldman2014parameterless}, P3 with
eLyMPuS~\citep{przewozniczek2026lympus}, the surrogate-assisted P3-GOMEA variants MO-DSA and MO-PSA of
Bartnik~\citep{bartnik2026evolutionary} (after~\citep{dushatskiy2021novelsurrogateassistedevolutionaryalgorithm}), and
multi-objective local search~\citep{denottelander2021localsearch}. \emph{General evolutionary multi-objective
algorithms:} NSGA-III~\citep{deb2014nsga3}, MOEA/D~\citep{zhang2007moead}, and SH-EMOA~\citep{guerreroviu2021bagofbaselines},
which builds on SMS-EMOA~\citep{beume2007smsemoa}. \emph{Bayesian optimisation:} MO-BANANAS~\citep{guerreroviu2021bagofbaselines,white2021bananas},
MOTPE~\citep{ozaki2022motpe} as implemented in Optuna~\citep{akiba2019optuna}, SMAC3 with
ParEGO~\citep{lindauer2022smac3,knowles2006parego}, and three Gaussian-process methods: OSS
Vizier~\citep{song2022ossvizier,song2024vizier}, and qNEHVI and qParEGO in BoTorch~\citep{daulton2021qnehvi,daulton2020qehvi,balandat2020botorch}.
\emph{Multi-fidelity:} Hyperband~\citep{li2018hyperband}, ASHA~\citep{li2020asha} with the multi-objective promotion
rule of MO-ASHA~\citep{schmucker2021moasha}, and the multi-fidelity versions of SH-EMOA, MO-BANANAS (MOSH-BANANAS),
and MO-BOHB~\citep{guerreroviu2021bagofbaselines,falkner2018bohb}. Random search~\citep{bergstra2012random} is the
reference point. Beyond JAHS-Bench-201 and NAS-HPO-Bench-II, the benchmarks are NAS-Bench-201 with energy on an
edge GPU from HW-NAS-Bench as the second objective~\citep{dong2020nasbench201,li2021hwnasbench} and the four FCNet
tabular tasks~\citep{klein2019tabular}. The experimental protocol follows the NAS best-practice
checklist of Lindauer and Hutter~\citep{lindauer2020bestpractices}.


% v005 stage C draft (2026-11-17): the v003 "LLM-based search" paragraph shortened to 4 sentences (v005-plan.md,
% stage C: "Akapit o LLM w Related Work skrócić do 3-4 zdań"). Keys as in v003; the v003 reason "why a
% linkage-learning engine fails to separate from random search" no longer holds and is dropped.

\paragraph{LLM-based search.} Large language models now act as the optimiser~\citep{zheng2023genius,nasir2024llmatic},
as a surrogate~\citep{hao2024llmsurrogate,jawahar2024llmpp}, or as a prior that warm-starts Bayesian
optimisation~\citep{liu2024llambo,zhang2024llmhpo}. Their documented gains concentrate in the early, sparse-data phase
and on small spaces, and several studies trace them to prior knowledge rather than to feedback from
evaluations~\citep{kristiadi2024sober,gupta2025llmbo,ferreira2026centaur}. Linkage learning is the opposite case: it
uses no prior and learns structure only from evaluations. We include no LLM-based arm, because it would draw on
information outside the evaluation budget that binds every other arm, and because NAS-Bench-201 is public and may be
part of the models' training data~\citep{zheng2023genius}.


\paragraph{Joint NAS+HPO benchmarks and their limits.} JAHS-Bench-201 answers queries through a learned surrogate
over a mixed categorical/continuous joint space~\citep{bansal2022jahsbench}; NAS-HPO-Bench-II is a fixed lookup
table over architecture, learning rate, and batch size~\citep{hirose2021nashpobenchii}. Only the latter admits an
exactly enumerable oracle Pareto front, which we exploit for IGD+ (Section~\ref{sec:results}). Guerrero-Viu et
al.\ establish SH-EMOA and MO-BOHB as baselines for this joint setting~\citep{guerreroviu2021bagofbaselines}, none
of them a linkage-learning EA. Both benchmarks share the NAS-Bench-201 cell space~\citep{dong2020nasbench201} and
small-scale image classification. The benchmarking literature documents that search-space design accounts for
much of a method's reported performance and that rankings transfer poorly across
benchmarks~\citep{yang2020nasfrustratinglyhard,mehta2022nasbenchsuite}, that this particular cell space is
redundant and skewed toward already-good architectures~\citep{lopes2023nasbenchdesign}, and that conclusions
transfer inconsistently across task domains~\citep{tu2022nasbench360,duan2021transnasbench101}. Any result on
this family should be read against that.
